% notes-substructures-and-quantifiers.tex
% Topic: Substructures and Quantifiers

\begin{remark*}[Exposition]
The preservation behavior of formulas depends on the quantifiers used in the
formula. Universal formulas tend to pass downward to substructures, while
existential formulas tend to pass upward to extensions. Mixed quantifier
patterns require more care.
\end{remark*}

\begin{definitionbox}{Definition ($\mathcal{L}$-Substructure)}
\begin{definition}[$\mathcal{L}$-Substructure]\label{def:l-substructure}
Let $\mathcal{A}$ and $\mathcal{B}$ be $\mathcal{L}$-structures.
The structure $\mathcal{A}$ is an \textbf{$\mathcal{L}$-substructure} of
$\mathcal{B}$, and $\mathcal{B}$ is an \textbf{extension} of $\mathcal{A}$, if
$A\subseteq B$ and each symbol of $\mathcal{L}$ is interpreted in
$\mathcal{A}$ as the restriction of its interpretation in $\mathcal{B}$.
We write $\mathcal{A}\subseteq\mathcal{B}$.
\end{definition}
\end{definitionbox}

\begin{remark*}[Standard quantified statement]
$\mathcal{A}\subseteq\mathcal{B}$ iff $A\subseteq B$,
$c^{\mathcal{A}}=c^{\mathcal{B}}$ for every constant symbol $c$,
$f^{\mathcal{A}}(\bar a)=f^{\mathcal{B}}(\bar a)$ for every function symbol
$f$ and tuple $\bar a\in A^{\operatorname{ar}(f)}$, and
\[
  \bar a\in R^{\mathcal{A}}
  \quad\text{iff}\quad
  \bar a\in R^{\mathcal{B}}
\]
for every relation symbol $R$ and tuple
$\bar a\in A^{\operatorname{ar}(R)}$.
\end{remark*}

\begin{remark*}[Interpretation]
Equivalently, $\mathcal{A}$ is a substructure of $\mathcal{B}$ exactly when
the inclusion map $A\hookrightarrow B$ is an $\mathcal{L}$-embedding.
\end{remark*}

\begin{dependencies}
\begin{itemize}
  \item \hyperref[def:l-structure]{$\mathcal{L}$-Structure}
  \item \hyperref[def:l-structure-embedding]{Embedding of $\mathcal{L}$-Structures}
\end{itemize}
\end{dependencies}

\subsection{Existential and Universal Formulas}

\begin{remark*}[Quantifier direction]
For $\mathcal{A}\subseteq\mathcal{B}$, universal assertions true in
$\mathcal{B}$ remain true in $\mathcal{A}$ because the elements of
$\mathcal{A}$ are among the elements quantified over in $\mathcal{B}$.
Existential assertions true in $\mathcal{A}$ remain true in $\mathcal{B}$
because witnesses from $A$ are still available in $B$.
\end{remark*}

\begin{example*}[Universal and existential behavior]
In the ordered language, the sentence
\[
  \forall x\,(0<x\lor x=0)
\]
is universal. If it holds in a structure, it also holds in any substructure.

By contrast, the sentence
\[
  \exists x\,(x+x=1)
\]
is existential. If it holds in a structure, it also holds in any extension,
because the same witness remains present in the larger universe.
\end{example*}

\begin{definitionbox}{Definition (Quantifier-Free Formula)}
\begin{definition}[Quantifier-Free Formula]\label{def:quantifier-free-formula}
A \textbf{quantifier-free formula} is a formula containing no quantifiers.
Equivalently, the quantifier-free formulas are defined recursively as follows:
\begin{enumerate}
  \item every atomic formula is quantifier-free;
  \item if $\varphi$ is quantifier-free, then $\lnot\varphi$ is
        quantifier-free;
  \item if $\varphi$ and $\psi$ are quantifier-free, then
        $(\varphi\land\psi)$ is quantifier-free.
  \item only formulas constructed by these rules in finitely many steps are
        quantifier-free.
\end{enumerate}
Other Boolean connectives are treated as abbreviations.
\end{definition}
\end{definitionbox}

\begin{remark*}[Standard quantified statement]
$\varphi$ is quantifier-free iff no occurrence of $\forall$ or $\exists$
appears in $\varphi$.
\end{remark*}

\begin{remark*}[Predicate reading]
\[
\mathsf{QuantifierFree}(\varphi)
\Longleftrightarrow
\lnot\mathsf{HasQuantifier}(\varphi).
\]
\end{remark*}

\begin{remark*}[Interpretation]
Quantifier-free formulas speak only about the currently assigned tuple. They
do not assert the existence of new witnesses or quantify over all elements of
the structure.
\end{remark*}

\begin{dependencies}
\begin{itemize}
  \item \hyperref[def:l-atomic-formula]{Atomic $\mathcal{L}$-Formula}
  \item \hyperref[def:l-formula]{$\mathcal{L}$-Formula}
\end{itemize}
\end{dependencies}

\begin{definitionbox}{Definition (Existential Formula)}
\begin{definition}[Existential Formula]\label{def:existential-formula}
An \textbf{existential formula} is defined recursively as follows:
\begin{enumerate}
  \item every quantifier-free formula is existential;
  \item if $\varphi$ is existential and $x$ is a variable, then
        $\exists x\,\varphi$ is existential;
  \item only formulas constructed by these rules in finitely many steps are
        existential.
\end{enumerate}
\end{definition}
\end{definitionbox}

\begin{remark*}[Standard quantified statement]
$\varphi$ is existential iff it is obtained from a quantifier-free formula by
prefixing finitely many existential quantifiers.
\end{remark*}

\begin{remark*}[Predicate reading]
\[
\mathsf{ExistentialFormula}(\varphi)
\Longleftrightarrow
\exists\psi\,\exists y_1,\ldots,y_n\,
\bigl(\mathsf{QuantifierFree}(\psi)
\land \varphi=\exists y_1\cdots\exists y_n\,\psi\bigr).
\]
\end{remark*}

\begin{remark*}[Interpretation]
Existential formulas assert that some finite tuple of witnesses exists making
a quantifier-free condition true.
\end{remark*}

\begin{dependencies}
\begin{itemize}
  \item \hyperref[def:quantifier-free-formula]{Quantifier-Free Formula}
  \item \hyperref[def:l-formula]{$\mathcal{L}$-Formula}
\end{itemize}
\end{dependencies}

\begin{definitionbox}{Definition (Universal Formula)}
\begin{definition}[Universal Formula]\label{def:universal-formula}
A \textbf{universal formula} is defined recursively as follows:
\begin{enumerate}
  \item every quantifier-free formula is universal;
  \item if $\varphi$ is universal and $x$ is a variable, then
        $\forall x\,\varphi$ is universal;
  \item only formulas constructed by these rules in finitely many steps are
        universal.
\end{enumerate}
\end{definition}
\end{definitionbox}

\begin{remark*}[Standard quantified statement]
$\varphi$ is universal iff it is obtained from a quantifier-free formula by
prefixing finitely many universal quantifiers.
\end{remark*}

\begin{remark*}[Predicate reading]
\[
\mathsf{UniversalFormula}(\varphi)
\Longleftrightarrow
\exists\psi\,\exists y_1,\ldots,y_n\,
\bigl(\mathsf{QuantifierFree}(\psi)
\land \varphi=\forall y_1\cdots\forall y_n\,\psi\bigr).
\]
\end{remark*}

\begin{remark*}[Interpretation]
Universal formulas assert that every finite tuple satisfies a quantifier-free
condition.
\end{remark*}

\begin{dependencies}
\begin{itemize}
  \item \hyperref[def:quantifier-free-formula]{Quantifier-Free Formula}
  \item \hyperref[def:l-formula]{$\mathcal{L}$-Formula}
\end{itemize}
\end{dependencies}

\subsection{Duality}

\begin{remark*}[Duality]
Every existential formula has the form
\[
  \exists y_1\cdots\exists y_n\,\psi
\]
with $\psi$ quantifier-free, and every universal formula has the form
\[
  \forall y_1\cdots\forall y_n\,\psi
\]
with $\psi$ quantifier-free. Since $\forall y$ may be treated as an
abbreviation for $\lnot\exists y\lnot$, a universal formula is equivalent to
the negation of an existential formula.
\end{remark*}

\subsection{Preservation Laws}

\begin{lemma}[Quantifier-Free Formulas Are Preserved by Substructures]\label{lem:quantifier-free-formulas-preserved-by-substructures}
Let $\mathcal{A}\subseteq\mathcal{B}$ be $\mathcal{L}$-structures. If
$\varphi(\bar x)$ is a quantifier-free $\mathcal{L}$-formula and
$\bar a\in A^n$, then
\[
  \mathcal{A}\models\varphi(\bar a)
  \quad\text{iff}\quad
  \mathcal{B}\models\varphi(\bar a).
\]
\hyperref[prf:quantifier-free-formulas-preserved-by-substructures]{\textit{Go to proof.}}
\end{lemma}

\begin{remark*}[Standard quantified statement]
For every substructure $\mathcal{A}\subseteq\mathcal{B}$, every
quantifier-free formula $\varphi(\bar x)$, and every tuple $\bar a\in A^n$,
truth of $\varphi(\bar a)$ is the same in $\mathcal{A}$ and $\mathcal{B}$.
\end{remark*}

\begin{remark*}[Predicate reading]
\[
\mathsf{Substructure}(\mathcal{A},\mathcal{B})
\land
\mathsf{QuantifierFree}(\varphi)
\Longrightarrow
\bigl(\mathcal{A}\models\varphi(\bar a)
\leftrightarrow
\mathcal{B}\models\varphi(\bar a)\bigr).
\]
\end{remark*}

\begin{remark*}[Interpretation]
Quantifier-free formulas only inspect the named relations, functions, and
constants on the given tuple. A substructure and its extension agree on that
local information.
\end{remark*}

\begin{dependencies}
\begin{itemize}
  \item \hyperref[def:l-substructure]{$\mathcal{L}$-Substructure}
  \item \hyperref[def:quantifier-free-formula]{Quantifier-Free Formula}
  \item \hyperref[lem:embeddings-preserve-atomic-formulas]{Embeddings Preserve Atomic Formulas}
\end{itemize}
\end{dependencies}

\begin{proposition}[Existential Formulas Are Preserved in Extensions]\label{prop:existential-formulas-preserved-in-extensions}
Let $\mathcal{A}\subseteq\mathcal{B}$ be $\mathcal{L}$-structures. Let
$\varphi(\bar x)$ be an existential $\mathcal{L}$-formula. If
$\bar a\in A^n$ and $\mathcal{A}\models\varphi(\bar a)$, then
$\mathcal{B}\models\varphi(\bar a)$.
\hyperref[prf:existential-formulas-preserved-in-extensions]{\textit{Go to proof.}}
\end{proposition}

\begin{remark*}[Standard quantified statement]
For every substructure $\mathcal{A}\subseteq\mathcal{B}$, every existential
formula $\varphi(\bar x)$, and every tuple $\bar a\in A^n$,
\[
  \mathcal{A}\models\varphi(\bar a)
  \Longrightarrow
  \mathcal{B}\models\varphi(\bar a).
\]
\end{remark*}

\begin{remark*}[Predicate reading]
\[
\mathsf{Substructure}(\mathcal{A},\mathcal{B})
\land
\mathsf{ExistentialFormula}(\varphi)
\Longrightarrow
\bigl(\mathcal{A}\models\varphi(\bar a)
\to
\mathcal{B}\models\varphi(\bar a)\bigr).
\]
\end{remark*}

\begin{remark*}[Interpretation]
Existential formulas are preserved upward because witnesses from the
substructure remain available in the extension.
\end{remark*}

\begin{dependencies}
\begin{itemize}
  \item \hyperref[def:l-substructure]{$\mathcal{L}$-Substructure}
  \item \hyperref[def:existential-formula]{Existential Formula}
  \item \hyperref[lem:quantifier-free-formulas-preserved-by-substructures]{Quantifier-Free Formulas Are Preserved by Substructures}
\end{itemize}
\end{dependencies}

\begin{proposition}[Universal Formulas Are Preserved in Substructures]\label{prop:universal-formulas-preserved-in-substructures}
Let $\mathcal{A}\subseteq\mathcal{B}$ be $\mathcal{L}$-structures. Let
$\varphi(\bar x)$ be a universal $\mathcal{L}$-formula. If
$\bar a\in A^n$ and $\mathcal{B}\models\varphi(\bar a)$, then
$\mathcal{A}\models\varphi(\bar a)$.
\hyperref[prf:universal-formulas-preserved-in-substructures]{\textit{Go to proof.}}
\end{proposition}

\begin{remark*}[Standard quantified statement]
For every substructure $\mathcal{A}\subseteq\mathcal{B}$, every universal
formula $\varphi(\bar x)$, and every tuple $\bar a\in A^n$,
\[
  \mathcal{B}\models\varphi(\bar a)
  \Longrightarrow
  \mathcal{A}\models\varphi(\bar a).
\]
\end{remark*}

\begin{remark*}[Predicate reading]
\[
\mathsf{Substructure}(\mathcal{A},\mathcal{B})
\land
\mathsf{UniversalFormula}(\varphi)
\Longrightarrow
\bigl(\mathcal{B}\models\varphi(\bar a)
\to
\mathcal{A}\models\varphi(\bar a)\bigr).
\]
\end{remark*}

\begin{remark*}[Interpretation]
Universal formulas are preserved downward because a substructure quantifies
over fewer elements than its extension.
\end{remark*}

\begin{dependencies}
\begin{itemize}
  \item \hyperref[def:l-substructure]{$\mathcal{L}$-Substructure}
  \item \hyperref[def:universal-formula]{Universal Formula}
  \item \hyperref[prop:existential-formulas-preserved-in-extensions]{Existential Formulas Are Preserved in Extensions}
\end{itemize}
\end{dependencies}
